% !PRINTED 261
\weberpage{261}

Ist $m$ ein Theiler von $p^2-1$, und

\[
\varrho = \gamma^{\frac{p^2-1}{m}},
\]
so ist $\varrho^m=1$ und es giebt keine niedrigere als die $m^{\mathrm{te}}$
Potenz von $\varrho$, die gleich 1 wird. Daher nennen wir $\varrho$ eine
primitive $m^{\mathrm{te}}$ Einheitswurzel (in $\mathfrak{E}_{2,p}$). Alle anderen $m^{\mathrm{ten}}$
Einheitswurzeln sind dann Potenzen $\varrho^s$ von $\varrho$ und unter diesen
sind die und nur die primitiv, bei denen $s$ relativ prim zu $m$ ist.

Dies folgt daraus, dass jede von Null verschiedene Zahl $\xi$
in der Form $\xi=\gamma^r$ darstellbar ist, und dass $\xi^m=\gamma^{rm}$ dann
und nur dann $=1$ ist, wenn $rm$ durch $p^2-1$ theilbar ist.

Jede von Null verschiedene Zahl in $\mathfrak{E}_{2,p}$ ist eine Einheitswurzel
 vom Grade $p^2-1$.

\begin{center}
\S.\,68.

\gesperrt{Die reelle lineare Congruenzgruppe $L_p$.}
\end{center}

Wir gehen nun, mit diesen Hülfsmitteln ausgestattet, an die
Untersuchung der reellen Congruenzgruppe $L_p$, die aus allen
Substitutionen der Form

\[
A = \begin{pmatrix} a, & b \\ c, & d \end{pmatrix}
\]
besteht, worin $a, b, c, d$ reelle, nach dem Modul $p$ genommene
ganze Zahlen sind, die der Bedingung $ad-bc\equiv 1$ genügen.
Die Gruppe $L_p$ ist nach der früheren Bezeichnung mit $E_{1,p}$ zu
bezeichnen, und ihr Grad ist, wenn wir den Fall $p=2$ ausschliessen:

\[
\frac{p(p^2-1)}{2}.
\]

Sind $A, U$ irgend zwei Elemente aus einer Gruppe, so haben
wir schon früher

\[
U^{-1}AU = A'
\]
das durch $U$ aus $A$ transformirte Element genannt. $A'^2, A'^3, \ldots$
sind transformirt aus $A^2, A^3, \ldots$, woraus folgt, dass alle aus
einander durch Transformation gewonnenen Elemente denselben
Grad haben. Entnehmen wir die Elemente $A$ aus irgend einer
Gruppe $G$, so bilden die transformirten Elemente eine mit $G$
isomorphe Gruppe

\[
U^{-1}GU = G',
\]
die wir die transformirte Gruppe von $G$ nennen.
